\documentclass[10pt,a4paper]{amsart}
\usepackage{amsmath,amssymb,amsthm,graphicx}
\usepackage{verbatim}
\newenvironment{reference}{\comment}{\endcomment}
%\newenvironment{reference}{}{}
\newenvironment{slogan}{\comment}{\endcomment}
\newenvironment{history}{\comment}{\endcomment}

% For commutative diagrams we use Xy-pic
\usepackage[all]{xy}

% We use 2cell for 2-commutative diagrams.
\xyoption{2cell}
% Theorem environments.
%
\theoremstyle{plain}
\newtheorem{theorem}[subsection]{Theorem}
\newtheorem{proposition}[subsection]{Proposition}
\newtheorem{lemma}[subsection]{Lemma}

\theoremstyle{definition}
\newtheorem{definition}[subsection]{Definition}
\newtheorem{example}[subsection]{Example}
\newtheorem{exercise}[subsection]{Exercise}
\newtheorem{situation}[subsection]{Situation}

\theoremstyle{remark}
\newtheorem{remark}[subsection]{Remark}
\newtheorem{remarks}[subsection]{Remarks}

\numberwithin{equation}{subsection}

% Macros
%
\def\lim{\mathop{\mathrm{lim}}\nolimits}
\def\colim{\mathop{\mathrm{colim}}\nolimits}
\def\Spec{\mathop{\mathrm{Spec}}}
\def\Hom{\mathop{\mathrm{Hom}}\nolimits}
\def\Ext{\mathop{\mathrm{Ext}}\nolimits}
\def\SheafHom{\mathop{\mathcal{H}\!\mathit{om}}\nolimits}
\def\SheafExt{\mathop{\mathcal{E}\!\mathit{xt}}\nolimits}
\def\Sch{\mathit{Sch}}
\def\Mor{\mathop{\mathrm{Mor}}\nolimits}
\def\Ob{\mathop{\mathrm{Ob}}\nolimits}
\def\Sh{\mathop{\mathit{Sh}}\nolimits}
\def\NL{\mathop{N\!L}\nolimits}
\def\CH{\mathop{\mathrm{CH}}\nolimits}
\def\proetale{{pro\text{-}\acute{e}tale}}
\def\etale{{\acute{e}tale}}
\def\QCoh{\mathit{QCoh}}
\def\Ker{\mathop{\mathrm{Ker}}}
\def\Im{\mathop{\mathrm{Im}}}
\def\Coker{\mathop{\mathrm{Coker}}}
\def\Coim{\mathop{\mathrm{Coim}}}

% Boxtimes
%
\DeclareMathSymbol{\boxtimes}{\mathbin}{AMSa}{"02}

%
% Macros for moduli stacks/spaces
%
\def\QCohstack{\mathcal{QC}\!\mathit{oh}}
\def\Cohstack{\mathcal{C}\!\mathit{oh}}
\def\Spacesstack{\mathcal{S}\!\mathit{paces}}
\def\Quotfunctor{\mathrm{Quot}}
\def\Hilbfunctor{\mathrm{Hilb}}
\def\Curvesstack{\mathcal{C}\!\mathit{urves}}
\def\Polarizedstack{\mathcal{P}\!\mathit{olarized}}
\def\Complexesstack{\mathcal{C}\!\mathit{omplexes}}
% \Pic is the operator that assigns to X its picard group, usage \Pic(X)
% \Picardstack_{X/B} denotes the Picard stack of X over B
% \Picardfunctor_{X/B} denotes the Picard functor of X over B
\def\Pic{\mathop{\mathrm{Pic}}\nolimits}
\def\Picardstack{\mathcal{P}\!\mathit{ic}}
\def\Picardfunctor{\mathrm{Pic}}
\def\Deformationcategory{\mathcal{D}\!\mathit{ef}}
\usepackage{cleveref}
\crefname{originaldefinition}{original definition}{original definitions}
\Crefname{originaldefinition}{Original definition}{Original definitions}
\makeatletter
\newlabel{definition-relative-dimension}{{00QD}{source}{}{sourcepointer-1}{}}
\newlabel{definition-relative-dimension@cref}{{[originaldefinition][0][]00QD}{[1][0][]source}{}{}{}}
\makeatother
\crefname{originallemma}{original lemma}{original lemmas}
\Crefname{originallemma}{Original lemma}{Original lemmas}
\makeatletter
\newlabel{lemma-CM-dim-finite-type}{{00R9}{source}{}{sourcepointer-2}{}}
\newlabel{lemma-CM-dim-finite-type@cref}{{[originallemma][0][]00R9}{[1][0][]source}{}{}{}}
\makeatother
\crefname{originallemma}{original lemma}{original lemmas}
\Crefname{originallemma}{Original lemma}{Original lemmas}
\makeatletter
\newlabel{lemma-dimension-at-a-point-finite-type-field}{{00P1}{source}{}{sourcepointer-3}{}}
\newlabel{lemma-dimension-at-a-point-finite-type-field@cref}{{[originallemma][0][]00P1}{[1][0][]source}{}{}{}}
\makeatother
\crefname{originallemma}{original lemma}{original lemmas}
\Crefname{originallemma}{Original lemma}{Original lemmas}
\makeatletter
\newlabel{lemma-dimension-fibres-bounded-open-upstairs}{{00QH}{source}{}{sourcepointer-4}{}}
\newlabel{lemma-dimension-fibres-bounded-open-upstairs@cref}{{[originallemma][0][]00QH}{[1][0][]source}{}{}{}}
\makeatother
\crefname{originallemma}{original lemma}{original lemmas}
\Crefname{originallemma}{Original lemma}{Original lemmas}
\makeatletter
\newlabel{lemma-one-equation}{{00KW}{source}{}{sourcepointer-5}{}}
\newlabel{lemma-one-equation@cref}{{[originallemma][0][]00KW}{[1][0][]source}{}{}{}}
\makeatother
\title{Being a regular sequence in the fibre is an open condition}
\author{The Stacks Project Authors}
\date{stacks-project:algebra.tex:local-snapshot-20261012; GNU FDL 1.2 or later}
\begin{document}
\maketitle
\noindent\textit{Editorial note.} All delivered mathematical text is the author's own, copied from the pinned Stacks Project chapter file algebra.tex; no statement or proof step was rewritten, reordered or removed.\par\medskip
\noindent\textit{Editorial note.} Preamble provenance, as required for this source: the frozen snapshot stages the chapter's own preamble file beside it (preamble.tex, 217 lines, 8116 bytes, the same shared Stacks-wide preamble the published chow and etale-cohomology cards use); the first line of the pinned chapter file is \textbackslash{}input\{preamble\}, so no preamble text lives inside the chapter itself. The card takes that file's lines 8-18 (the verbatim package and the three \textbackslash{}comment-based environments reference/slogan/history at lines 9-12, section 13.2, together with the Xy-pic setup) and its lines 154 to the runtime end of the file (the native theorem declarations and mathematical macros). Not one character was changed, no package was added, no macro was defined, and the excerpt serves only as preamble, never as a span coordinate.\par\medskip
\noindent\textit{Editorial note.} The same-chapter pointers are located at their \textbackslash{}label lines inside the pinned chapter file; the cross-chapter pointers are located in the project's own tag list staged in the snapshot (sources/stacks-chow/tags.txt), at the verbatim lines carrying their canonical tags.\par\medskip
\noindent\textit{Editorial note.} The counted proof refers to the definition of relative dimension and to same-chapter statements; the definition is recorded with kind definition in the navigation stubs, and the arguments stay inside the counted span.\par\medskip
\noindent\textit{Editorial note.} The counted spans contain no \textbackslash{}cite, so the card carries no bibliography.\par\medskip
\noindent\textit{Allowed standard prerequisites (editorial).}\begin{itemize}
\item Relative dimension at a point. Used in the counted proof at lines 33498-33505. Supplies the numerical invariant that is compared with the length of the sequence.
\item Cohen-Macaulayness read through the dimension function. Invoked at the end of the counted proof (lines 33518-33520). Converts the numerical inequality into regularity of the sequence.
\end{itemize}\medskip
\ifdefined\hypertarget\hypertarget{sourcepointer-1}{}\fi
\noindent\textit{Source reference.} Original definition 00QD: sources/stacks-algebra/pinned-original-source.tex, source lines 31367--31367. Navigation pointer for \textbackslash{}ref\{definition-relative-dimension\} used inside the counted span (line 33499); the argument that consumes it is retained in the span.\par\medskip
\ifdefined\hypertarget\hypertarget{sourcepointer-2}{}\fi
\noindent\textit{Source reference.} Original lemma 00R9: sources/stacks-algebra/pinned-original-source.tex, source lines 33453--33453. Navigation pointer for \textbackslash{}ref\{lemma-CM-dim-finite-type\} used inside the counted span (line 33520); the argument that consumes it is retained in the span.\par\medskip
\ifdefined\hypertarget\hypertarget{sourcepointer-3}{}\fi
\noindent\textit{Source reference.} Original lemma 00P1: sources/stacks-algebra/pinned-original-source.tex, source lines 28439--28439. Navigation pointer for \textbackslash{}ref\{lemma-dimension-at-a-point-finite-type-field\} used inside the counted span (line 33503); the argument that consumes it is retained in the span.\par\medskip
\ifdefined\hypertarget\hypertarget{sourcepointer-4}{}\fi
\noindent\textit{Source reference.} Original lemma 00QH: sources/stacks-algebra/pinned-original-source.tex, source lines 31525--31525. Navigation pointer for \textbackslash{}ref\{lemma-dimension-fibres-bounded-open-upstairs\} used inside the counted span (line 33513); the argument that consumes it is retained in the span.\par\medskip
\ifdefined\hypertarget\hypertarget{sourcepointer-5}{}\fi
\noindent\textit{Source reference.} Original lemma 00KW: sources/stacks-algebra/pinned-original-source.tex, source lines 14911--14911. Navigation pointer for \textbackslash{}ref\{lemma-one-equation\} used inside the counted span (line 33506); the argument that consumes it is retained in the span.\par\medskip
\begin{lemma}
\label{lemma-open-regular-sequence}
Let $R \to S$ be a finite type ring map. Let $d$ be an integer such that all
fibres $S \otimes_R \kappa(\mathfrak p)$ are Cohen-Macaulay and equidimensional
of dimension $d$. Let $f_1, \ldots, f_i$ be elements of $S$. The set
$$
\{ \mathfrak q \in V(f_1, \ldots, f_i)
\mid f_1, \ldots, f_i
\text{ are a regular sequence in }
S_{\mathfrak q}/\mathfrak p S_{\mathfrak q}
\text{ where }\mathfrak p = R \cap \mathfrak q
\}
$$
is open in $V(f_1, \ldots, f_i)$.
\end{lemma}

\begin{proof}
Write $\overline{S} = S/(f_1, \ldots, f_i)$.
Suppose $\mathfrak q$ is an element of the set defined in the
lemma, and $\mathfrak p$ is the corresponding prime of $R$.
We will use relative dimension as defined in
Definition \ref{definition-relative-dimension}.
First, note that $d = \dim_{\mathfrak q}(S/R) =
\dim(S_{\mathfrak q}/\mathfrak pS_{\mathfrak q}) +
\text{trdeg}_{\kappa(\mathfrak p)}\ \kappa(\mathfrak q)$
by Lemma \ref{lemma-dimension-at-a-point-finite-type-field}.
Since $f_1, \ldots, f_i$ form a regular sequence in the
Noetherian local ring $S_{\mathfrak q}/\mathfrak pS_{\mathfrak q}$
Lemma \ref{lemma-one-equation} tells us that
$\dim(\overline{S}_{\mathfrak q}/\mathfrak p\overline{S}_{\mathfrak q})
= \dim(S_{\mathfrak q}/\mathfrak pS_{\mathfrak q}) - i$.
We conclude that $\dim_{\mathfrak q}(\overline{S}/R)
= \dim(\overline{S}_{\mathfrak q}/\mathfrak p\overline{S}_{\mathfrak q}) +
\text{trdeg}_{\kappa(\mathfrak p)}\ \kappa(\mathfrak q)
= d - i$ by Lemma \ref{lemma-dimension-at-a-point-finite-type-field}.
By Lemma \ref{lemma-dimension-fibres-bounded-open-upstairs}
we have $\dim_{\mathfrak q'}(\overline{S}/R) \leq d - i$
for all $\mathfrak q' \in V(f_1, \ldots, f_i) = \Spec(\overline{S})$
in a neighbourhood of $\mathfrak q$. Thus after replacing
$S$ by $S_g$ for some $g \in S$, $g \not \in \mathfrak q$
we may assume that the inequality holds for all
$\mathfrak q'$. The result follows from Lemma
\ref{lemma-CM-dim-finite-type}.
\end{proof}

\end{document}
